% Target language: Traditional Chinese (zh-Hant).
% Source commit: 0d064f6ce5c6f02a2e89edc6004b56723e8042be.
% Status: AI-assisted Traditional Chinese localization; native-speaker review requested.
\StoryParagraph{3}{一顆卵子的重量，三點五微克。}

\StoryParagraph{3}{五千萬美元現金的重量，兩點五噸。}

\StoryParagraph{3}{景甜在蒙太奇拉古納海灘的電話裡向我要的是後者，抵押的是前者。}

\StoryParagraph{3}{她去那裡是為了取卵。代孕的孩子預計2027年出生，快一點的話還能屬馬。這是她主動提的。到了診所，她說沒有五千萬，不取。}

\UnprintedPause

\StoryParagraph{0}{蒙太奇拉古納海灘在海崖上，底下是太平洋。那一層包下來，三十天，一百萬美元，為了她的隱私。她住了八天就走了。}

\StoryParagraph{0}{剩下的二十二天，酒店的清潔工每天照常上去。是我的助理講的，她去交鑰匙的時候撞見過一次：一個墨西哥女人推著車，一間一間進，換毛巾，擦檯面，把被角掖好，退出來把門帶上。做完一整層，兩個小時。}

\StoryParagraph{0}{助理問她，沒人住，為什麼還要做。}

\StoryParagraph{0}{她說，登記上寫著有人。}

\StoryParagraph{0}{一座失去皇后的凡爾賽宮。}

\StoryParagraph{0}{我突然意識到，景甜喜歡這個。景甜喜歡一個地方因為她而空著。什麼都不發生。}

\StoryParagraph{3}{與蒙太奇拉古納海灘一起空置的，還有一架灣流G700尷尬地停在Van Nuys，機組每天照常報到。}

\StoryParagraph{3}{第八天她打了那個電話。之後沒有人再見過她。}

\NarrativeBreak

\StoryParagraph{0}{2007年，已經有人替她燒錢了。那年我剛進北大。}

\StoryParagraph{3}{當然，燒錢的不是我。}

\StoryParagraph{3}{QQ彈出一個視窗，代言人是她。就在那臺聯想IBM的老機器上，解析度糊得很。我把圖片存下來了。}

\StoryParagraph{3}{我在一年房租一千塊的北大宿舍上網，在校內網上找到她，發了好友申請。校內網要填一句留言，我寫了刪，刪了寫，弄了大概四十分鐘，最後發出去一句“你好”。}

\StoryParagraph{3}{她沒有接受。}

\StoryParagraph{3}{她也許根本沒用過校內網。}

\StoryParagraph{3}{後來這些年，我做每一件事的時候都在想一個問題：到了什麼程度，她才會接受。第一個十億的時候我想過，一百億的時候我想過。}

\StoryParagraph{3}{後來我不想了。人人網早就倒閉了。}

\StoryParagraph{3}{我沒有再申請第二次。}

\StoryParagraph{3}{2011年，北京冷。我買了一件羽絨服，一百五十塊。我記得這個價格，因為我猶豫了三天。}

\StoryParagraph{3}{同一年她拍了一部一億五千萬的電影《戰國》。孫紅雷、吳鎮宇、金喜善、中井貴一、姜武就像十九年後的蒙太奇拉古納海灘與灣流G700一樣被請來做了陪襯，他們不知道自己為什麼出現在這裡，來都來了，總得發生些什麼，但是什麼都沒有發生。}

\StoryParagraph{3}{再後來是一億五千萬美元的《長城》。張藝謀，劉德華，馬特·達蒙。再後來是成龍。景甜宇宙的名單越來越長。}

\StoryParagraph{3}{我那時已經在美國了，讀書，創業，越來越忙。唯一沒變的是，每換一臺裝置，無論多麻煩，我都把那張QQ圖片導過去。}

\StoryParagraph{3}{十九年後，名單上加了我。}

\NarrativeBreak

\StoryParagraph{3}{瑰麗酒店我們見了第一面。我一直以為見到她那天我會緊張。}

\StoryParagraph{3}{我沒有。}

\StoryParagraph{3}{她比照片上瘦，穿得很隨意，靠在沙發上。我看著她，忽然想不起來我等的到底是什麼。}

\StoryParagraph{3}{我帶她去看卡特蘭的香蕉。她看了很久，問我，這個多少錢。}

\StoryParagraph{3}{六百二十萬美元。}

\StoryParagraph{3}{她說，那它爛了怎麼辦。}

\StoryParagraph{3}{我沉默了。她說，不管了，我們去看電影吧。}

\StoryParagraph{3}{我說好。}

\StoryParagraph{3}{她說放映廳裡除了我們兩個，不能有任何人。}

\StoryParagraph{3}{我說好。}

\StoryParagraph{3}{我出去打了個電話。回來的時候她還站在那根香蕉前面。}

\StoryParagraph{3}{那部電影是《瘋狂動物城2》。}

\StoryParagraph{3}{兩個人的放映廳，這種場景一般用於中宣部審片。}

\StoryParagraph{3}{我讓助理把電影院所有的座位都買下來。}

\StoryParagraph{3}{我看過批判刷假票房的文章，說票買光了只有兩個人看。現在我明白了，也許不全是為了票房。}

\StoryParagraph{3}{把座位全買下來不難，麻煩的是已經賣出去的票。助理和保鏢提前到影院門口，遇到買了票的觀眾，把現金遞過去。}

\StoryParagraph{3}{一對帶孩子的夫妻來了，孩子已經在門口高高舉著爆米花。}

\StoryParagraph{3}{助理遞上一沓錢，孩子不懂為什麼不能進去，大人接過錢，道謝，什麼也沒問，走了。}

\StoryParagraph{3}{那天領錢的有二十六個人。}

\StoryParagraph{3}{看完她把微信頭像換成了朱迪警官。}

\StoryParagraph{3}{晚上我用手揉她的奶子，被她推開了。}

\StoryParagraph{3}{我以為她不願意。我把手收回來，準備說點什麼。}

\StoryParagraph{3}{她沒說話。她坐起來，從包裡翻出一個小袋子，拉開，裡面一排東西，像一套很小的工具。}

\StoryParagraph{3}{她說，你的指甲扎得我痛。}

\StoryParagraph{3}{我看了一眼。我一直是拿指甲刀隨便剪兩下，邊上是尖的，平時不覺得。}

\StoryParagraph{3}{她說，過來。}

\StoryParagraph{3}{她把我的手放在她腿上，一根一根磨。}

\StoryParagraph{3}{先是粗的那一面，磨完換一面，再換一面。磨一會兒舉起來對著燈看看，再接著磨。}

\StoryParagraph{3}{我說這個要磨多久。}

\StoryParagraph{3}{她說，磨到不刮人為止。}

\StoryParagraph{3}{十個手指，她花了快一個小時。}

\StoryParagraph{3}{中間她教我，這一面磨形狀，這一面拋光，最後一面不要來回蹭，要順著一個方向。}

\StoryParagraph{3}{我說好。}

\StoryParagraph{3}{磨完她把我的手翻過來看了看，摸了一下。}

\StoryParagraph{3}{她說，行了。}

\StoryParagraph{3}{摸吧，想摸哪都行。}

\StoryParagraph{3}{中間她說過一句話。她說，要是以後我不在了，就再沒有人給你拋指甲了。}

\StoryParagraph{3}{我笑了一下，沒說話。}

\StoryParagraph{0}{認識她三個月的時候，我讓人做了一份東西。她那些年的專案、合作方、公司股權、景國慶，田心愛，所有的幾個關聯帳戶，北京租的房，上海買的房，還有一些別的。一週就做完了，一共四十幾頁。還有些男男女女的事情，我覺得這是最不重要的。}

\StoryParagraph{0}{裡面有一些事她後來跟我講過，講的和上面寫的不太一樣。我沒有說。}

\StoryParagraph{0}{有一些事她一直沒講。我也沒有問。}

\StoryParagraph{0}{我全部相信她說的。}

\StoryParagraph{0}{那份東西我留著，一直留到現在。}

\StoryParagraph{3}{但是我一個字也不信。}

\NarrativeBreak

\StoryParagraph{3}{她說，香港人太多，會被認出來，我們出去走走。}

\StoryParagraph{3}{我說好。}

\StoryParagraph{3}{她在地圖上指了一下特納裡費。}

\StoryParagraph{3}{我說好。}

\StoryParagraph{3}{我出去打了個電話，回來才知道那個地方位於非洲西北海岸外的大西洋上。第二天中午，我們一行三十個人出發了，A330在機場等著。}

\StoryParagraph{3}{飛機的臥室裡她問我，我的名字好聽嗎。}

\StoryParagraph{3}{我說好聽。}

\StoryParagraph{3}{她說，我爸爸的姓和我媽媽的姓拼在一起的，是不是很巧。}

\StoryParagraph{3}{我說我第一天就記住了。}

\StoryParagraph{3}{她抱著我說，以後別叫我景甜了，叫我媽媽吧。}

\StoryParagraph{3}{好，媽媽。}

\StoryParagraph{3}{落地以後機長過來告訴我，老闆，島上一開始不讓降，他們沒見過這麼大的私人飛機。}

\StoryParagraph{3}{島上規矩一樣，周圍不能有人。助理去處理了，沒有人問為什麼，這次我沒問多少人。}

\StoryParagraph{3}{麗茲卡爾頓正對著大西洋，落地窗太大。她睡覺不能有一絲光。助理用黑膠帶把整個房間貼起來，外面亮如白晝，裡面黑得像禁閉室。膠帶一曬就掉，每天要重貼。}

\StoryParagraph{3}{助理一共十一個人，早上六點半到，等我們睡了才走。貼膠帶的是兩個女孩，一個負責窗框，一個負責空調和插座的指示燈，貼完再檢查三遍。膠帶買了三十卷，後來又補了二十卷。她們發現下午三點以後貼的最容易掉，因為那時候玻璃最燙，後來改成早上貼。這件事沒有人教她們，是她們自己試出來的。她們互相之間會討論，像討論一項工藝。}

\StoryParagraph{3}{她知道那兩個女孩叫什麼。我不知道。}

\StoryParagraph{3}{我們在泰德公園的公路邊停下看日落，助理從車上搬下來玫瑰花，開了香檳。}

\StoryParagraph{3}{媽媽看著雲海說，好美。不過不如青海德令哈的日落。}

\StoryParagraph{3}{我是青海人。}

\StoryParagraph{3}{我沒有說。}

\StoryParagraph{3}{夜晚我們在山頂看星星，穿著棉襖還是冷。}

\StoryParagraph{3}{媽媽說，許個願吧。}

\StoryParagraph{3}{我問，你許了什麼。}

\StoryParagraph{3}{媽媽說，不告訴你。}

\StoryParagraph{3}{媽媽說，太美了。但就是太冷了，我們去馬爾地夫好不好。}

\StoryParagraph{3}{我說，好吧。}

\StoryParagraph{3}{特納裡費零下十五度，馬爾地夫二十七度，直線距離九千五百公里，飛十一個小時。不過馬爾地夫的機場太小了，我們寬體客機中途得換成小飛機。}

\StoryParagraph{3}{媽媽說，小飛機聲音好吵。還是大飛機好。}

\StoryParagraph{3}{去馬爾地夫的飛機臥室裡，媽媽跟我說，她許的願是，和我要一個孩子。必須代孕。}

\StoryParagraph{3}{我沒有問過為什麼。我以為她怕疼。}

\StoryParagraph{3}{中途媽媽突然坐起來找手機，連上星鏈，說糟了。}

\StoryParagraph{3}{我問怎麼了。}

\StoryParagraph{3}{她說，有個人要威脅我，問了三天了，我一直沒回。}

\StoryParagraph{3}{我說要多少。}

\StoryParagraph{3}{她說不用你管。}

\StoryParagraph{3}{她打字打了很久，刪了改，改了刪。打完把手機扣在腿上，過一會兒又翻過來看。}

\StoryParagraph{3}{我問那人是誰。}

\StoryParagraph{0}{她說，以前跟我一起的，財務，進去了，扛了事。}

\StoryParagraph{3}{出來走投無路，一定要找我。}

\StoryParagraph{3}{我說，乾脆報警抓了。}

\StoryParagraph{3}{她看了我一眼。你不懂。}

\StoryParagraph{3}{我說你給了嗎。}

\StoryParagraph{3}{她說給了。}

\StoryParagraph{3}{多少。}

\StoryParagraph{3}{八萬。}

\StoryParagraph{3}{我說這個也用得著想三天。}

\StoryParagraph{3}{她沒說話。}

\StoryParagraph{3}{那天剩下的時間她一直沒怎麼說話。落地了才睡著。}

\StoryParagraph{0}{最早到索尼娃賈尼的是媽媽的箱子。三十個人的行李，小飛機裝不下，分了三架。媽媽只有一隻箱子，單獨一架，她說不要和別人的放在一起。}

\StoryParagraph{0}{索尼娃賈尼是出了名的no news, no shoes，媽媽放心了。}

\StoryParagraph{3}{酒店管家，保鏢，助理，一共配了二十六個人，還有兩個無人機手遠遠跟著。助理們又開始貼膠帶了，空調的燈，插座的燈，一遍又一遍。}

\StoryParagraph{0}{晚上我們像交換暗號一樣交換了彼此的前任，就像一個王朝在核對年號。媽媽不願意有任何差池，不過媽媽有個禁忌詞，那個運動員不能提。}

\StoryParagraph{0}{我沒提，雖然我知道，媽媽背上的紋身是因為運動員紋的。}

\StoryParagraph{3}{可能因為太大了洗不掉。}

\StoryParagraph{0}{她說，現在我們是一家人了。我爸媽把我養這麼大不容易，你得給彩禮。}

\StoryParagraph{0}{我說好。}

\StoryParagraph{0}{她報了兩個名字，景國慶、田心愛，和兩個帳號。}

\StoryParagraph{0}{三千萬，人民幣。}

\StoryParagraph{0}{我轉了。那兩個名字，我是在那四十幾頁裡第一次見到的。}

\StoryParagraph{3}{過了一會兒，她把手機遞過來給我看。田心愛回了兩個字：收到。}

\StoryParagraph{0}{有一次媽媽說，你知道我要這些東西的時候在想什麼嗎。}

\StoryParagraph{0}{我說什麼。}

\StoryParagraph{0}{她說，我在想，你會不會有一次說不。}

\StoryParagraph{0}{我說不會。}

\StoryParagraph{0}{她說，那就沒意思了。}

\StoryParagraph{3}{我想，我又有什麼做不到嗎？}

\StoryParagraph{0}{太無聊的時候，我在擺弄那張2007年的QQ圖片，想傳到新裝置上。媽媽把頭湊過來，你怎麼揹著我存別的女人的照片，這是誰。}

\StoryParagraph{0}{我沉默了。什麼也沒說。}

\StoryParagraph{0}{她看圖太糊了，就去忙別的了。}

\StoryParagraph{0}{那天晚上我把那張圖刪了。}

\StoryParagraph{3}{第二天又從舊手機裡找回來了。}

\StoryParagraph{0}{媽媽手機也刷個不停，我問她看什麼。}

\StoryParagraph{0}{她說沒什麼，就是撕番位的那些破事。}

\StoryParagraph{0}{撕番位，一部戲只能有一個明星排第一。男一女一也只能有一個人排第一。}

\StoryParagraph{0}{我說無聊啊。}

\StoryParagraph{0}{她說，是啊。更難的是，有時候劇組夫妻，一進組四五十天，一般和一個男的拍戲，晚上做愛，白天還要撕番位。}

\StoryParagraph{0}{我說太反人性了。}

\StoryParagraph{0}{她說是啊。}

\StoryParagraph{0}{不過拍完戲了，無論之前做過多少次愛，就當一切都沒發生。}

\StoryParagraph{3}{然後她又看了半個小時。}

\StoryParagraph{0}{2026年1月2日，索尼娃賈尼，新的一年，我向媽媽求婚了。}

\StoryParagraph{0}{是1月1日早上決定的。戒指在香港。島上沒有賣戒指的地方，助理坐水上飛機去了馬累，下午回來，戒指裝在一個酒店的信封裡。}

\StoryParagraph{0}{媽媽說，鑽戒太小了。}

\StoryParagraph{0}{媽媽說，沒事，以後補個大的。}

\StoryParagraph{0}{另外，你能來北京嗎？}

\StoryParagraph{0}{我沉默了。}

\StoryParagraph{0}{她以為我不想去。}

\StoryParagraph{0}{媽媽說，你為什麼對我這麼好。}

\StoryParagraph{0}{我說因為我喜歡你。}

\StoryParagraph{0}{她說，那你以前對別人也這樣嗎。}

\StoryParagraph{0}{我說沒有。}

\StoryParagraph{3}{她笑了一下，沒說話。}

\StoryParagraph{0}{有一天早上我醒得早，她不在床上。}

\StoryParagraph{0}{我出去看，她坐在露臺上，對著海，整個海是黑的。沒開燈，天還沒亮。}

\StoryParagraph{0}{我問她怎麼起這麼早。}

\StoryParagraph{0}{她說，我在算日子。}

\StoryParagraph{0}{我說算什麼。}

\StoryParagraph{0}{她說，你回去睡吧。}

\StoryParagraph{0}{我在她旁邊坐了一會兒，不知道說什麼。後來我回去睡了。}

\StoryParagraph{3}{我再醒過來的時候她已經化好妝了。}

\StoryParagraph{0}{在馬爾地夫待了十四天，沒有一個人認出她。}

\StoryParagraph{0}{第十三天晚上在餐廳，她把墨鏡摘了。過了一會兒，帽子也摘了，放在桌上。她就那樣坐了一頓飯。隔壁桌是一對德國老夫妻，服務生是斯里蘭卡人，沒有人多看她一眼。}

\StoryParagraph{0}{吃完她把帽子戴回去了。}

\StoryParagraph{0}{媽媽說，我們回去吧。}

\StoryParagraph{0}{我說不是還沒住夠嗎。}

\StoryParagraph{3}{她說，太安靜了。}

\UnprintedPause

\StoryParagraph{0}{回到香港，媽媽見了我父母。}

\StoryParagraph{0}{當天晚上她說，家裡要籌備兩家見面和婚事，儘可能多打點。還是景國慶和田心愛。}

\StoryParagraph{0}{我沒問為什麼。}

\StoryParagraph{0}{我說一天一張卡限額一百萬。}

\StoryParagraph{0}{媽媽說太慢了。}

\StoryParagraph{0}{轉五張。}

\StoryParagraph{0}{五張卡，兩個帳號，輪著轉。轉到第三張的時候，我猶豫了一下。我還是轉完了。}

\StoryParagraph{3}{媽媽說，效率太低了。}

\StoryParagraph{0}{我說，過年把咱爹媽接到香港來一起過吧。}

\StoryParagraph{0}{媽媽說，今年已經定好在海南過年了。明年吧。}

\StoryParagraph{0}{春節，維港放煙花，來了一百萬人。沒有她。}

\StoryParagraph{0}{凌晨三點，我把父母送回去。}

\StoryParagraph{0}{路上很空，紅燈沒有一輛車在等，我還是停了。}

\StoryParagraph{0}{媽媽說先不見了，馬上去洛杉磯取卵了。手頭的事情要先處理完。}

\StoryParagraph{0}{媽媽又問，能來北京嗎？}

\StoryParagraph{0}{我沉默了。}

\StoryParagraph{3}{她沒有說北京有什麼。我也沒問。}

\StoryParagraph{0}{媽媽說，香港必須買一套房子，寫她的名字。}

\StoryParagraph{0}{我說，好吧。}

\StoryParagraph{3}{情人節也沒見面。}

\StoryParagraph{0}{有一次她在電話裡說，今天早上抽了八管血。}

\StoryParagraph{0}{我說這麼多。}

\StoryParagraph{0}{她說還好。}

\StoryParagraph{0}{我說辛苦了。}

\StoryParagraph{0}{然後我們聊了別的。}

\StoryParagraph{0}{後來我才知道，那種抽血要空腹，要在月經週期的第二天或者第三天，日子不能錯。她那天早上大概六點就出門了，北京的冬天，天還沒亮。後來的檢查越來越多。我知道這件事，是因為助理每月的彙總付款單裡多了一項。第一次出現是去年，金額不大，幾千塊。後來每個月都有，有時候兩三筆。我批得很快。這類金額不需要我看明細。後來才知道那些是抽血、激素六項、B超、AMH。要看她還有多少儲備，能取出來幾個，值不值得往下走。}

\StoryParagraph{0}{這些都是在北京做的。}

\StoryParagraph{3}{我不知道是哪家醫院。我也不知道她是自己去的，還是有人陪。}

\StoryParagraph{0}{2月25日她落地香港。}

\StoryParagraph{0}{她說不去洛杉磯了，除非先看到香港的房子。}

\StoryParagraph{0}{我把過去半年看的十幾套房都發給她。她都嫌不好看。那十幾套房我一套也沒去過，是助理花了半年一間一間挑的，助理轉給我，我轉給她。有一套在半山，影片裡能看見維港，拍的是傍晚。有一個瞬間鏡頭晃了一下，拍到了拍影片的人的影子在牆上。我看著那個影子想，這個人是誰，他今天走了多少路，看了多少房，他知不知道這套房是給誰看的。}

\StoryParagraph{3}{時間太緊，她還是上了去洛杉磯的G700，2月28日飛到蒙太奇拉古納海灘。}

\NarrativeBreak

\StoryParagraph{0}{第八天媽媽打了電話。她那邊是下午，我這邊是凌晨。}

\StoryParagraph{0}{她說，五千萬，美元。}

\StoryParagraph{0}{我說，讓我想想。}

\StoryParagraph{0}{她沒有說話。我聽見她那邊有海。}

\StoryParagraph{0}{過了一會兒她說，好。}

\StoryParagraph{0}{她掛了。}

\StoryParagraph{3}{我認識她以來，沒有對她說過“讓我想想”。}

\StoryParagraph{0}{我用Claude Code讀API把所有現金資產跑了一遍，結論是對我不會有任何影響。我又核了一遍。我在Claude Code裡輸入了問題，Claude說，不要把這五千萬美元給她。}

\StoryParagraph{3}{我問，她不再愛我了嗎？}

\StoryParagraph{3}{Claude說，對於愛情，我不關心，也不理解。但是你不能把這五千萬美元給她。}

\StoryParagraph{3}{我從來沒有拒絕過她，不是慷慨，是害怕。害怕下一秒的答案。}

\StoryParagraph{0}{而現在有人替我說了。}

\StoryParagraph{0}{我盯著那行字看了很久。讓我毛骨悚然的不是它聰明，是Claude一點猶豫都沒有。}

\StoryParagraph{0}{它不為難。}

\StoryParagraph{0}{我當時想，這大概是更高的一種東西。人做不到這樣。}

\StoryParagraph{3}{她不再愛我了嗎。}

\StoryParagraph{3}{第二天她打回來。}

\StoryParagraph{3}{媽媽說，海很大。}

\StoryParagraph{3}{我說是嗎。}

\StoryParagraph{3}{她說，這層樓太大了，有回聲。早上有人來打掃，我躲在衛生間裡等她們走。}

\StoryParagraph{3}{我說可以不讓她們來。}

\StoryParagraph{3}{她說不用。}

\StoryParagraph{3}{她說，診所問我，孩子的父親什麼時候到。}

\StoryParagraph{3}{我說我可以派人過去。}

\StoryParagraph{3}{她說，不用。}

\StoryParagraph{3}{然後她說，想好了嗎。}

\StoryParagraph{3}{我沒有說話。不是在想，我沒有在想任何事。電話裡她那邊有海，我這邊是維港，兩邊都沒有人說話。我數了一下，十一秒。}

\StoryParagraph{3}{景甜說，行。}

\StoryParagraph{3}{景甜說，我知道了。}

\StoryParagraph{3}{她掛了。}

\StoryParagraph{3}{我後來無數次回想過，那兩天她除了五千萬還說了什麼。海很大，這層樓太大，她躲在衛生間裡。就這些，我全記得，一個字不差。我記不起來的是她的聲音——她說“我知道了”的時候，是什麼語氣。}

\StoryParagraph{0}{我想不起來。後來我就不想了。}

\StoryParagraph{0}{我把手機放在床頭櫃上，螢幕朝下。然後我做了一件事，至今不知道為什麼：我把手放在胸口，數了一下。}

\StoryParagraph{0}{心跳很穩。}

\StoryParagraph{0}{我又數了一遍，還是很穩。一分鐘六十幾下，和平時沒有區別。我坐在那裡，手按著自己的胸口，像一個醫生在給別人做檢查。}

\StoryParagraph{3}{我想，原來是這樣。}

\StoryParagraph{0}{那天剩下的時間我什麼也沒幹。我在房間裡待到天黑又到天亮，看著維多利亞港從藍色變成灰色再變成黑色。早晨六點，維港對面的金融中心一棟一棟地開始亮燈。中間助理敲過一次門，問晚餐怎麼安排，我說不用了。}

\StoryParagraph{3}{我沒有哭。我一直在等我哭。}

\NarrativeBreak

\StoryParagraph{0}{後來助理跟我結賬，裡面還有一項是特納裡費的酒店賠償。我們走了以後，膠帶又貼了幾天，沒人告訴她們不用貼了。酒店說貼得太久，次數太多，窗框上的膠痕清不掉，那面玻璃要整塊換掉。}

\StoryParagraph{0}{我說批了。}

\StoryParagraph{3}{那面玻璃朝著大西洋。}

\StoryParagraph{0}{帳單裡飛機燃油那一項，比蒙太奇那一層樓還貴。我問了一下。助理說，A330是按兩百人設計的，我們只有三十個人，太輕了，起飛的配平不對。所以每次都要多加油，加到夠重為止。到了地方，多加的那些又不能留著，得放掉。}

\StoryParagraph{0}{我說，每次都這樣？}

\StoryParagraph{3}{他說，每次都這樣。}

\StoryParagraph{0}{還有洛杉磯診所寄回來的一個箱子。八天的促排針，要冷藏的，一支沒拆。}

\StoryParagraph{3}{還有代孕機構的定金，二十萬美元，不退。代孕的人一月就定好了，加州人，三十四歲，生過兩個。檔案裡有她的照片，我沒看過。她一月就準備好了。}

\StoryParagraph{3}{景甜走後，我把那張2007年的照片翻出來看了一次。十九年了，裝置畫素越來越高，而這張圖，彷彿畫素越來越低，臉都是糊的。我放大看了很久，想找出她和後來有哪裡不一樣。一處也沒找出來。因為那張照片本來就看不清。}

\StoryParagraph{0}{有一天我用指甲刀剪指甲，剪完摸了一下，是尖的。}

\StoryParagraph{0}{我想起來抽屜裡還有那個東西，她留下的。}

\StoryParagraph{0}{我拿出來磨了兩下，不對。她說過要順著一個方向，我忘了是哪個方向。}

\StoryParagraph{0}{我把它放回去了。}

\StoryParagraph{3}{我的指甲現在還是磨人的。}

\StoryParagraph{0}{有一天早上手機彈出一個提醒，是助理去年排日程的時候設的：預產期前三個月，準備房間。}

\StoryParagraph{3}{我把它劃掉了。}

\StoryParagraph{0}{她走了以後，我一直在想一個問題。}

\StoryParagraph{0}{如果那天我給了，她會留下來嗎。}

\StoryParagraph{0}{我想了很久。}

\StoryParagraph{3}{後來我把那份東西拿出來，又看了一遍。四十幾頁。}

\StoryParagraph{0}{然後我想起春節那天晚上，在頂樓俯瞰香港維港的煙花。}

\StoryParagraph{0}{一百萬人像潮水一樣湧來，煙花結束，一百萬人退去。}

\StoryParagraph{0}{一百萬人。}

\StoryParagraph{0}{每年都是這樣。}

\StoryParagraph{0}{凌晨三點的香港，}

\StoryParagraph{3}{什麼都沒有發生。}
